\documentclass{article}

\usepackage{amsmath, amsthm, amssymb, amsfonts}
\usepackage{thmtools}
\usepackage{graphicx}
\usepackage{setspace}
\usepackage{geometry}
\usepackage{float}
\usepackage{hyperref}
\usepackage[utf8]{inputenc}
\usepackage[english]{babel}
\usepackage{framed}
\usepackage[dvipsnames]{xcolor}
\usepackage{tcolorbox}
\usepackage{physics}

\DeclareMathOperator{\grd}{grad}
\DeclareMathOperator{\crl}{curl}
\DeclareMathOperator{\dvr}{div}
\DeclareMathOperator{\length}{length}
\colorlet{LightGray}{White!90!Periwinkle}
\colorlet{LightOrange}{Orange!15}
\colorlet{LightGreen}{Green!15}

\newcommand{\HRule}[1]{\rule{\linewidth}{#1}}

\declaretheoremstyle[name=Theorem,]{thmsty}
\declaretheorem[style=thmsty,numberwithin=section]{theorem}
\tcolorboxenvironment{theorem}{colback=LightGray}

\declaretheoremstyle[name=Proposition,]{prosty}
\declaretheorem[style=prosty,numberlike=theorem]{proposition}
\tcolorboxenvironment{proposition}{colback=LightOrange}

\declaretheoremstyle[name=Principle,]{prcpsty}
\declaretheorem[style=prcpsty,numberlike=theorem]{principle}
\tcolorboxenvironment{principle}{colback=LightGreen}

\setstretch{1.2}
\geometry{
    textheight=9in,
    textwidth=6in,
    top=0.5in,
    headheight=12pt,
    headsep=25pt,
    footskip=30pt
}

% ------------------------------------------------------------------------------

\begin{document}


% ------------------------'

\textbf{Final Exam}

\textbf{Exercise 1.} In this exercise, you will derive a formula for the distance between
two lines $L_1(t) = A_1t + B_1$ and $L_2(t) = A_2t + B_2$ in $\mathbb{R}^3$ (that is, the minimal distance between any pair $P,Q$ of points where $P$ lies on $L_1$ and $Q$ on $L_2$).

(a) Let $f: X \to \mathbb{R}$ be a function on some set $X$. Let $Y \subseteq X$
and suppose $y \in Y$ is a global minimum for $f$. That is, 
\[f(y) \leq f(x)\]
for all $x \in X$. Show/explain why $y$ is also a global minimum for
$f|_Y: Y \to \mathbb{R}$, where $f|_A$ is the restriction of $f$ to $Y$. 
(Don't overthink this. This is simply a helpful observation for later.)

(b) Assume for simplicity that $L_1$ and $L_2$ are not parallel. Find equations
(in terms of the $A_i$ and $B_i$) for two parallel planes, one containing $L_1$ and the other containing $L_2$.
In particular, note that there exist such planes. (Hint: there is a very natural expression for the normal vector.)

(c) Recall that the \emph{projection} of a vector $A$ onto a vector $B$ is 
\[\frac{A \cdot B}{B \cdot B} B.\]
Use this to find the distance between the two planes (this is the length of any perpendicular segment connecting the two planes). Your answer should
be an expression involving $A_1$, $A_2$, $B_1$, and $B_2$. 

(d) Observe that if two non-parallel lines $L_1$ and $L_2$ lie in parallel planes (as above), then
there is a point $P$ on $L_1$ and a point $Q$ on $L_2$ such that the segment $PQ$ 
is perpendicular to both planes. In light of part (a), explain why
the distance you found in part (c) is indeed the distance between the lines $L_1$ and $L_2$.

\newpage
\textbf{Exercise 2.} If we drag a line segment around in the plane, we sweep out
some area (think of a roller paint brush). Similarly, if we drag a planar figure (such as a square or disk) around in 3-space,
we obtain some kind of solid. The volume $V$ of this solid is given by Guldin's formula, which states
\[V = \int_{t_0}^{t_1} A(t) \mathbf{n}\cdot C'(t) dt,\]
where $C(t)$ is the curve traced out by the centroid/center of mass of the planar figure, $\mathbf{n}$ is the unit normal of the planar figure,
and $A(t)$ is the area of the planar figure. Note that $A(t)$ can depend on time, meaning
the figure is allowed to change shape as it traverses through space. Note also that $\mathbf{n}\cdot C'(t)$ is
simply the component of the velocity of the centroid of the figure along the direction perpendicular to the figure.

(a) Use this formula to derive the volume of a cone with a base of radius $r$ and
a height $h$. (Hint: a cone is obtained by dragging a disk upward at constant speed, while letting
the radius shrink to $0$ linearly with time. If we center the cone around the $z$-axis, the trajectory of the 
centroid is simply $C(t) = (0,0,t)$, where $0 \leq t \leq h$).

(b) Now suppose that we skew the cone by letting the cross section $A$ ``drift'' in the $x$ and $y$
directions, so that the centroid now traces out $C(t) = (k_1 t, k_2 t, t)$. What is the volume of the
resulting solid?

\newpage
\textbf{Exercise 3.} Recall that for any two vectors $A$ and $B$,
\[|A \cdot B| \leq \norm{A} \norm{B}.\] 
Also recall that 
\[\left| \int_a^b f(x)dx  \right| \leq \int_a^b |f(x)|dx.\]
We also have seen that if $f(x) \leq g(x)$ for all $x \in [a,b]$, then \[\int_a^b f(x)dx \leq \int_a^b g(x)dx.\]
Recall once again that the length of a curve is the integral of speed:
\[\int_{a}^{b} \norm{C'(t)} dt.\]
Suppose now that $F: \mathbb{R}^3 \to \mathbb{R}^3$ is a vector field that is \textbf{bounded}. This means there is some
number $M$ such that
$|F(X)| < M$ for all $X \in \mathbb{R}^3$. Use the above facts to show that for a curve $C(t)$ defined 
for $a \leq t \leq b$, we have
\[\int_a^b F(C(t))\cdot C'(t)dt \leq M \length(C).\]
In other words, on a bounded vector field, the path integral $\int F \cdot dC$ will have a small value
if the path is short.

\newpage
\textbf{Exercise 4.} Let $C$ be the curve of intersection of the cylinder 
$x^2+y^2 = 2y$ and the plane $y=z$. Use Stokes' theorem to show that
\[\int_C y^2\ dx + xy\ dy + xz\ dz = 0.\]
(Hint: think about the relation between the normal vector of the plane $y=z$ and the curl of the field $F=(y^2,xy,dz)$.)

\end{document}